\documentclass[11pt,a4paper]{article}

\usepackage[margin=2.5cm]{geometry}
\usepackage{amsmath,amssymb,amsthm}
\usepackage{booktabs}
\usepackage{array}

\newtheorem{theorem}{Theorem}[section]
\newtheorem{conjecture}{Conjecture}[section]

\newcommand{\R}{\mathbb{R}}
\newcommand{\Z}{\mathbb{Z}}
\newcommand{\Q}{\mathbb{Q}}

\title{\textbf{A Dictionary for Matching Coefficients}}
\author{Robin Langer and Claude Code}
\date{August 2026}

\begin{document}
\maketitle

\begin{abstract}
\noindent
The matching coefficients of an orthogonal-polynomial family --- the
coefficients approximating the weighted functions $W^{1/2}\widehat{P}_m$
by the polynomials $\widehat{P}_k$ --- have been computed for ten
families in a series of standalone papers. Each computation, taken
alone, deflates to a classical mechanism. What does not deflate is the
three-column language in which the computations become entries: a
dictionary between analysis (generating-function kernels), combinatorics
(transfer-matrix automata) and arithmetic (generator lattices). This
paper records the dictionary, the four claims only it can state, and the
one family --- Wilson --- on which it makes a falsifiable prediction.
\end{abstract}

%% ===================================================================
\section{The Dictionary}
%% ===================================================================

The Askey scheme contains almost no theorems and is among the most-used
objects in the field. This paper is an object of the same kind.

\begin{center}
\small
\begin{tabular}{lll}
\toprule
\textbf{analysis} & \textbf{combinatorics} & \textbf{arithmetic} \\
\midrule
GF kernel & automaton & generator lattice \\
ODE order of the kernel & number of states & lattice rank \\
Taylor coefficients $u_n$ at base point & walk weights & all signs and generators \\
Cayley/Laguerre substitution & species chassis & universal, sign-free \\
kernel singularity type & --- & obstruction level \\
D-module reducibility & --- & algebraic vs.\ transcendental generators \\
\bottomrule
\end{tabular}
\end{center}

Four sentences can be said in this language and, so far as we know, in
no classical one.

\begin{enumerate}
\item \emph{Rank classifies.} The Gaussian lattice has rank $1$;
Meixner--Pollaczek has rank $2$ over Bessel (resp.\ Tricomi) values;
continuous Hahn has rank $2$ with one \emph{algebraic} generator, the
Jordan--H\"older filtration of its reducible D-module visible in the
lattice; Wilson admits no finite rank at all. No classical reference
makes this classification.
\item \emph{The combinatorics is universal; the family chooses only
$u_n$.} The defect--chain--matching chassis is identical across
families; the weight enters through the prefactor exponent and the
Taylor coefficients $u_n$ of the kernel, which carry every sign.
Proved for MP at $\lambda=\tfrac14$ and, by changing two lines of the
proof, at $\lambda=\tfrac34$.
\item \emph{Degeneration is evaluation, not limit.} The monic Jacobi
matching entries are rational functions of the parameter $\alpha$; the
Chebyshev closed forms are their values at $\alpha=-\tfrac12$, not an
analytic limit.
\item \emph{Obstructions form a hierarchy.} Entire weight $\to$
geometric decay and integrality over $\Z$; branch point at finite
distance $\to$ power-law decay, lattice survives; Barnes factor
$\Gamma(2ix)$ $\to$ no lattice.
\end{enumerate}

%% ===================================================================
\section{The Chassis}
%% ===================================================================

Fix a weight $W$ on $\R$ with orthonormal polynomials $\widehat{P}_n$
and set $\Phi_m = W^{1/2}\widehat{P}_m$ (for the Gamma-type weights, the
analytic split $\Gamma(\lambda+ix)\widehat{P}_m$ replaces the square
root). The \emph{forward approximation} coefficients are
$c_{mk} = \langle \Phi_m, \widehat{P}_k\rangle_W$; the \emph{reverse
expansion} coefficients are
$\langle \widehat{P}_m, \widehat{P}_k \rangle_{W^{1/2}}$. Both are
organised by a master generating function
\[
  \sum_{m,k} c_{mk}\, w^m u^k
  = \text{prefactor}(w,u)\; g(\zeta), \qquad \zeta = \varphi(w)\varphi(u),
\]
with $\varphi$ a Cayley or Laguerre substitution. The substitution
factors as $\zeta = (1+\sigma)(1+\tau)$ in rescaled variables, and
extracting coefficients produces one combinatorial chassis for every
family: defect sets from the prefactor (double-factorial weights),
chains from $\sigma,\tau$ (Lah numbers), and a Charlier-type partial
matching of chains. The kernel $g$ contributes only its Taylor
coefficients at the base point $\zeta=1$.

When $g$ satisfies a linear ODE of order $r$ with polynomial
coefficients, differentiating the ODE at the base point yields an
integer recurrence for the rescaled coefficients $u_n$, hence
$u_n \in \Z A_1 \oplus \cdots \oplus \Z A_r$ for the $r$ initial
values $A_i$. This is the whole arithmetic mechanism. The worked
instance is MP at $\lambda=\tfrac14$, where
$g = e^{\zeta/2}K_0(\zeta/2) = \sqrt{\pi}\,U(\tfrac12,1,\zeta)$ gives
$u_{n+2} = -2n\,u_{n+1} + (4n+2)\,u_n$ with $u_0 = A$, $u_1 = A-B$
\cite{mp14}; at $\lambda=\tfrac34$ the same two steps give
$u'_{n+2} = -2n\,u'_{n+1} + (4n+6)\,u'_n$ over the Tricomi generators
\cite{mp34}. Legendre's formula for the $2$-adic valuation of $n!$
clears every denominator; the walk weights are positive integers and
every sign of every table lives in $u_n$.

%% ===================================================================
\section{The Entries}
%% ===================================================================

One paragraph per family; proofs and tables are in the standalone
papers.

\begin{center}
\footnotesize
\begin{tabular}{lllll}
\toprule
family & GF kernel & order & lattice / generators & status \\
\midrule
Hermite & $\exp(\text{quadratic})$ & 1 & $\Z$; signed bipartite matchings & proved \cite{herm} \\
Laguerre & $2/(1+w+u-3wu)$ & 1 & $\Z$; Delannoy numbers & proved \cite{lag} \\
Charlier & Charlier kernel & 1 & $\Z$; one-state automaton & proved \cite{char} \\
Chebyshev--Jacobi & Wallis/Beta mass & --- & quartic classes over $\Gamma(\tfrac14)^2/\sqrt{2\pi}$ & proved \cite{jac} \\
MP $\lambda=\tfrac14$, forward & $e^{\zeta/2}K_0(\zeta/2)$ & 2 & $\Z A \oplus \Z B$, $A,B = \sqrt{2e}\,K_{0,1}(\tfrac12)$ & proved \cite{mp14} \\
MP $\lambda=\tfrac14$, reverse & $e^{-1/\zeta}$ & 1 & $(2\pi/e)\Z$ & proved \cite{mp14} \\
MP $\lambda=\tfrac34$ & $U(\tfrac32,1,\zeta)$ & 2 & $\Z A' \oplus \Z B'$, $A',B' = U(\tfrac32,1,1), U(\tfrac52,2,1)$ & proved \cite{mp34} \\
unbalanced $\Gamma_{1/4}\Gamma_{3/4}$ & first-order & 1 & single $A'' = e\,E_1(1)$ & proved \cite{unb} \\
continuous Hahn & $\sqrt{\zeta}\;U(1,1,\cdot)$-type & 2 & $\Q\mathcal{A} \oplus \Q\mathcal{B}$, $\mathcal{A} = \sqrt{\pi}e^2E_1(2)$, $\mathcal{B} = \sqrt{\pi}$ & empirical \cite{hahn} \\
Wilson & Barnes $\Gamma(2ix)$ & $\ge 3$? & none found to $10^{18}$ & open \cite{wil} \\
\bottomrule
\end{tabular}
\end{center}

Three of the entries carry the dictionary's weight. The MP pair is the
proved instance of claim~2. Continuous Hahn is the instance of the
reducibility row, its kernel built on $e^{z}E_1(z)$, whose D-module is
an extension of the trivial module by a Kummer module; the lattice sees
the filtration through the algebraic generator $\sqrt{\pi}$, where the
irreducible MP module has two transcendental generators. Chebyshev is
the instance of claim~3, recovered from Jacobi by evaluating rational
functions of $\alpha$ at $-\tfrac12$ \cite{jac}.

The Weil explicit formula provides a remark in the same language. In
the MP $\lambda=\tfrac14$ basis, the archimedean side compresses to the
rank-$2$ lattice $\Q\sqrt{2\pi} \oplus \Q\sqrt{2\pi}(\gamma+\log 8\pi)$
(proved), while the prime side shows no $\Q$-compression among $16$
entries (empirical) --- the explicit formula's two halves are
algebraically incommensurable in this basis \cite{weil}.

%% ===================================================================
\section{The Obstruction Hierarchy}
%% ===================================================================

The decay of $c_{0k}$ and the arithmetic degrade together, and the
kernel's singularity type is the common cause.

\begin{enumerate}
\item \emph{Entire weight} (Hermite, Laguerre, Charlier). Geometric
decay; integrality over $\Z$; closed combinatorial formulas.
\item \emph{Branch point at finite distance} (MP, Hahn; the branch
point of $\Gamma(\lambda+ix)^2\Gamma(\lambda-ix)$-type kernels at
$i/4$). Power-law decay of the coefficients (exponent $\approx
\tfrac34$ at $\lambda=\tfrac14$); the lattice survives, with rank the
ODE order.
\item \emph{Barnes factor} $\Gamma(2ix)$ (Wilson). No lattice found;
see \S\ref{sec:test}.
\end{enumerate}

The dual-pair construction repairs arithmetic, not analysis: replacing
the orthogonal basis by the biorthogonal pair
$\Gamma(\lambda\pm ix)\widehat{P}_m$ restores integrality at
$\lambda=\tfrac14$ without improving the $L^2$ truncation rate.

%% ===================================================================
\section{The Test}\label{sec:test}
%% ===================================================================

The dictionary predicts that lattice rank equals the ODE order of the
GF kernel. Wilson is the falsifiable case.

\begin{conjecture}\label{conj:wilson}
The generating-function kernel of the Wilson matching coefficients
satisfies no linear ODE of order $\le 2$ with polynomial coefficients,
and correspondingly the coefficients generate a $\Q$-vector space of
unbounded dimension as $m,k \to \infty$.
\end{conjecture}

Evidence. The entries $c_{mk}$, $m,k \le 5$, are $\Q$-combinations of
$11$ moments; PSLQ at $200$ digits with coefficient bound $10^{18}$
finds no relation among the moments beyond the exact
polynomial-product identities, so their $\Q$-rank is $\ge 16$ counting
both parity sectors --- against MP and Hahn, which compress to $2$
generators with coefficients below $10^{8}$. What is missing is the
analytic half: a computation of the D-module rank of the Wilson kernel
itself. If that rank were finite and small, the dictionary would be
wrong, and the failure would be as informative as the confirmations.
If the language cannot make the Wilson question well-posed --- the
kernel row of the table is presently conjectural for Wilson --- that
too should be recorded, and this section is the record.

%% ===================================================================
\section{Status of the Claims}
%% ===================================================================

\begin{center}
\begin{tabular}{lll}
\toprule
claim & status & source \\
\midrule
Hermite, Laguerre, Charlier integrality & proved & \cite{herm,lag,char} \\
Chebyshev--Jacobi classes; connection-Gram & proved & \cite{jac} \\
MP $\tfrac14$ forward $\Z A \oplus \Z B$ & proved & \cite{mp14} \\
MP $\tfrac14$ reverse $(2\pi/e)\Z$ & proved & \cite{mp14} \\
MP $\tfrac34$ $\Z A' \oplus \Z B'$ & proved & \cite{mp34} \\
unbalanced, single generator & proved & \cite{unb} \\
Hahn rank-$2$ lattice & empirical (PSLQ) & \cite{hahn} \\
Hahn generator $\sqrt{\pi}$ from reducibility & structural reading & \cite{hahn} \\
rank $=$ ODE order, in general & predicted & --- \\
Wilson, no finite lattice & empirical (to $10^{18}$) & \cite{wil} \\
Weil archimedean rank $2$ & proved & \cite{weil} \\
Weil prime side rank $\ge 16$ & empirical & \cite{weil} \\
\bottomrule
\end{tabular}
\end{center}

Empirical entries are PSLQ identifications at $\ge 40$ digits,
confirmed at doubled precision with coefficient bound $\ge 10^{10}$;
no empirical claim is used in the proof of any proved claim.

\begin{thebibliography}{9}
\bibitem{synth} R.~Langer and Claude Code,
\emph{Approximating Orthogonal Functions with Orthogonal Polynomials},
preprint, 2026.
\bibitem{herm} \emph{The Hermite Matching Matrix}, preprint, 2026.
\bibitem{lag} \emph{The Laguerre Matching Matrix}, preprint, 2026.
\bibitem{char} \emph{The Charlier Matching Matrix}, preprint, 2026.
\bibitem{jac} \emph{The Jacobi Matching Matrix}, preprint, 2026.
\bibitem{mp14} \emph{The Dual Meixner--Pollaczek Matching Coefficients:
the Calculation}, preprint, 2026.
\bibitem{mp34} \emph{The Dual Meixner--Pollaczek Matching Coefficients
at $\lambda = \tfrac34$}, preprint, 2026.
\bibitem{unb} \emph{The Unbalanced Dual Pair}, preprint, 2026.
\bibitem{hahn} \emph{Continuous Hahn Matching Coefficients}, preprint,
2026.
\bibitem{wil} \emph{Wilson Matching Coefficients}, preprint, 2026.
\bibitem{weil} \emph{The Weil Form in the Meixner--Pollaczek Basis},
preprint, 2026.
\end{thebibliography}

\end{document}
